\chapter{Implémentation C++ et ingénierie logicielle}
\label{chap:implementation}

\epigraph{\textit{``Any fool can write code that a computer can understand. Good programmers write code that humans can understand.''}}{--- Martin Fowler}

\section{Organisation du dépôt}

Le projet est organisé selon une architecture modulaire séparant clairement les interfaces, les implémentations, et les points d'entrée.

\subsection{Structure des répertoires}

\begin{lstlisting}[language=bash, caption={Arborescence du projet WFC}]
Projet801/
├── include/wfc/         # Headers publics
│   ├── Grid.hpp           # Domaine 2D
│   ├── Tile.hpp           # Sous-bloc N×N + helpers symétries
│   ├── Bitset.hpp         # Bitset packé sur uint64_t
│   ├── TileSet.hpp        # Catalogue de tuiles + fréquences
│   ├── OverlapRules.hpp   # Compatibilités tuile×offset
│   ├── Wave.hpp           # Superposition par cellule
│   ├── WFCSolver.hpp      # Interface solver + SolverOptions/Stats
│   ├── GridIO.hpp         # I/O txt, PPM, PNG
│   ├── solvers/           # Backends spécialisés
│   │   ├── WFCSolverSerial.hpp
│   │   ├── WFCSolverOMP.hpp
│   │   └── WFCSolverKokkos.hpp
│   └── internal/          # Helpers et base class
│       ├── WFCSolverBase.hpp
│       └── SolverCommon.hpp
├── src/                 # Implémentations correspondantes
├── apps/                # Points d'entrée
│   ├── wfc_serial.cpp
│   ├── wfc_omp.cpp
│   ├── wfc_kokkos.cpp
│   ├── benchmark.cpp
│   └── wfc_dungeon.cpp    # Démo UE5 (opt-in BUILD_DUNGEON)
├── tests/               # Suite de tests unitaires (15 suites)
├── samples/             # Grilles d'entrée
├── scripts/             # SLURM, plot, analyze, render
├── results/             # CSV de benchmarks
├── docs/                # Documentation projet (ARCHITECTURE, ALGORITHM, etc.)
├── ue5_plugin/          # Plugin Unreal Engine 5.7
└── third_party/         # stb_image_write.h
\end{lstlisting}

\subsection{Modules et responsabilités}

Le tableau~\ref{tab:modules} détaille la responsabilité de chaque module.

\begin{table}[H]
\centering
\begin{tabularx}{\textwidth}{lX}
\toprule
\textbf{Module} & \textbf{Responsabilité} \\
\midrule
\texttt{Grid.hpp} & Domaine 2D row-major en \texttt{uint8\_t}, accès toroïdal via \texttt{at\_torus(r, c)} \\
\midrule
\texttt{Tile.hpp} & Pattern $N \times N$ POD-like avec hash FNV-1a et helpers de symétrie (\texttt{rotated\_90}, \texttt{reflected\_horizontal}) \\
\midrule
\texttt{Bitset.hpp} & Trois classes : \texttt{Bitset} propriétaire, \texttt{BitsetView} mutable, \texttt{ConstBitsetView} lecture-seule. Opérations 64-bit packées. \\
\midrule
\texttt{TileSet.hpp/cpp} & Extraction toroïdale $N \times N$, déduplication, expansion D4 optionnelle \\
\midrule
\texttt{OverlapRules.hpp/cpp} & Construction parallèle (OMP) des règles de compatibilité \\
\midrule
\texttt{Wave.hpp} & Buffer plat de bitsets, NUMA first-touch, accès toroïdal pour BFS \\
\midrule
\texttt{WFCSolverBase.hpp/cpp} & Boucle solve, dispatch séquentiel/parallèle/backtrack, gestion stats \\
\midrule
\texttt{SolverCommon.hpp/cpp} & Helpers partagés : \texttt{weighted\_entropy}, \texttt{weighted\_pick}, \texttt{cell\_jitter}, \texttt{serial\_propagate}, \texttt{serial\_run\_attempt}, \texttt{serial\_run\_attempt\_backtrack} \\
\midrule
\texttt{WFCSolverSerial.cpp} & $\sim 25$ lignes : \texttt{pick\_cell} = \texttt{serial\_min\_entropy}, \texttt{propagate} = \texttt{serial\_propagate} \\
\midrule
\texttt{WFCSolverOMP.cpp} & Backend OpenMP : \texttt{parallel\_min\_entropy}, \texttt{propagate\_tasks}, \texttt{atomic\_and\_with} \\
\midrule
\texttt{WFCSolverKokkos.cpp} & Backend Kokkos GPU-portable : \texttt{Kokkos::View}, \texttt{PropagateOp} fonctor, \texttt{kHostOnly} compile-time branch \\
\bottomrule
\end{tabularx}
\caption{Modules du projet et leurs responsabilités.}
\label{tab:modules}
\end{table}

% =============================================================================
\section{Structures de données critiques}
\label{sec:impl:structures}
% =============================================================================

\subsection{Bitset packé sur uint64\_t}

\begin{important}{Pourquoi pas \texttt{std::bitset<N>} ou \texttt{std::vector<bool>} ?}
Trois raisons :
\begin{enumerate}
    \item \textbf{Taille dynamique requise} : \texttt{std::bitset<N>} exige $N$ en \texttt{constexpr}, incompatible avec un $L$ déterminé à runtime selon l'échantillon.
    \item \textbf{Vectorisation} : \texttt{std::vector<bool>} interdit l'accès direct aux mots sous-jacents. On veut faire des \texttt{\_\_builtin\_popcountll} et des AND/OR 64 bits à la fois. Notre \texttt{Bitset} expose \texttt{words()} directement.
    \item \textbf{Allocation} : \texttt{std::vector<bool>} reste un wrapper sur \texttt{vector<u8>}, pas plus dense. Notre \texttt{Bitset} est exactement $\lceil L / 64 \rceil$ mots de 64 bits.
\end{enumerate}
\end{important}

Pour $L = 11$ (l'exemple du sujet), le bitset tient dans un seul mot. Pour $L = 73$ (terrain $N = 3$), dans deux mots. Toujours en cache L1.

\subsection{Wave à buffer plat}

\begin{important}{Pourquoi pas \texttt{vector<Bitset>} ?}
Naïvement on mettrait un \texttt{vector<Bitset>} (un Bitset par cellule). Mauvais :
\begin{itemize}
    \item Chaque \texttt{Bitset} possède sa propre \texttt{vector<u64>}, autant d'allocations que de cellules. Pour $256 \times 256$, ce sont 65 536 allocations.
    \item Indirection mémoire supplémentaire à chaque accès cellule (pointeur vers le buffer interne du vector).
\end{itemize}
\end{important}

Choix : un seul \texttt{vector<u64>} contigu de taille $r \cdot c \cdot \text{words\_per\_cell}$, et les accesseurs \texttt{at(r, c)} retournent une \texttt{BitsetView} non-propriétaire qui pointe dans ce buffer. Une seule allocation pour toute la wave, layout cache-friendly, NUMA-friendly.

\begin{figure}[H]
\centering
\schemaImage{Layout flat de la Wave avec BitsetView}
\caption{Wave en mémoire : un seul buffer contigu, partagé par tous les threads, avec accès via BitsetView.}
\label{fig:wave_layout}
\end{figure}

\subsection{First-touch parallèle pour NUMA}

Le constructeur de \texttt{Wave} initialise les cellules en parallèle :

\begin{lstlisting}[language=C++]
#pragma omp parallel for schedule(static)
for (std::ptrdiff_t c = 0; c < total_cells; ++c) {
    const std::size_t base = c * words_per_cell_;
    for (std::size_t w = 0; w + 1 < words_per_cell_; ++w) {
        words_[base + w] = 0xFFFFFFFFFFFFFFFFULL;  // tous bits set
    }
    if (words_per_cell_ > 0) {
        words_[base + (words_per_cell_ - 1)] = tail_mask;
    }
}
\end{lstlisting}

Avec \texttt{OMP\_PROC\_BIND=close} et \texttt{OMP\_PLACES=cores}, le thread qui touche en premier la mémoire est aussi celui qui la consommera ensuite. Sur EPYC 9654 (8 NUMA nodes), cela évite que le thread 0 first-touche tout et que 7/8 des accès traversent un boundary NUMA.

\subsection{OverlapRules à plat}

Les règles d'adjacence sont stockées dans un \texttt{vector<Bitset>} indexé par
\[
\text{idx} = t \cdot (2N-1)^2 + \text{offset\_index}(dx, dy)
\]
avec $\text{offset\_index}(dx, dy) = (dx + N - 1) \cdot (2N - 1) + (dy + N - 1)$.

Pour le sample du sujet ($L = 11$, $N = 2$), cela représente $11 \times 9 \times 1 \times 8 = 792$ octets de table. Tient en L1, accédée intensivement pendant la propagation.

% =============================================================================
\section{Conception orientée objet : Template Method}
\label{sec:impl:tm}
% =============================================================================

L'architecture des solvers utilise le pattern \textbf{Template Method} :

\begin{lstlisting}[language=C++]
class WFCSolverBase : public WFCSolver {
public:
    // La méthode de plus haut niveau est finale.
    Grid solve(const TileSet&, const OverlapRules&,
               const SolverOptions&, SolverStats&) override final;

protected:
    // Hot path : à surcharger par chaque backend.
    virtual int pick_cell(...) = 0;
    virtual bool propagate(...) = 0;

    // Hooks optionnels (Kokkos init/finalize).
    virtual void on_solve_begin() {}
    virtual void on_solve_end() {}
    virtual const char* backend_name() const = 0;

private:
    // Boucle séquentielle, dispatch parallel-attempts/backtrack.
    Grid solve_sequential(...);
    Grid solve_parallel(...);
};
\end{lstlisting}

\textbf{Avantages :}
\begin{itemize}
    \item Le retry, la mesure de stats, l'orchestration parallel-attempts, le dispatch backtrack sont écrits \textbf{une seule fois} dans la classe de base.
    \item Chaque backend (\texttt{Serial}, \texttt{OMP}, \texttt{Kokkos}) ne fournit que les deux primitives \texttt{pick\_cell} et \texttt{propagate} qui le distinguent.
    \item Les invariants de l'orchestration (déterminisme, ordre des seeds) sont garantis identiquement pour tous les backends.
\end{itemize}

% =============================================================================
\section{Build : trois libs CMake}
\label{sec:impl:cmake}
% =============================================================================

Le \texttt{CMakeLists.txt} produit trois libs distinctes plutôt qu'une lib monolithique avec \texttt{ifdef} :

\begin{table}[H]
\centering
\begin{tabularx}{\textwidth}{lXl}
\toprule
\textbf{Cible CMake} & \textbf{Sources} & \textbf{Dépend de} \\
\midrule
\texttt{wfc\_core} & Grid, TileSet, OverlapRules, Wave, GridIO, SolverCommon, WFCSolverBase, WFCSolverSerial & OpenMP::OpenMP\_CXX si \texttt{USE\_OMP=ON} (pour la construction parallèle des règles) \\
\midrule
\texttt{wfc\_omp\_lib} & WFCSolverOMP & wfc\_core, OpenMP::OpenMP\_CXX \\
\midrule
\texttt{wfc\_kokkos\_lib} & WFCSolverKokkos & wfc\_core, Kokkos::kokkos \\
\bottomrule
\end{tabularx}
\caption{Cibles CMake et leurs dépendances.}
\label{tab:cmake_targets}
\end{table}

\textbf{Bénéfices :} \texttt{wfc\_serial} (binaire) ne lie pas OpenMP s'il n'est pas demandé. \texttt{wfc\_omp} ne lie pas Kokkos. La dépendance \texttt{Kokkos::kokkos} reste localisée à \texttt{wfc\_kokkos\_lib}, invisible pour qui n'utilise pas Kokkos.

\subsection{Options CMake}

\begin{table}[H]
\centering
\begin{tabularx}{\textwidth}{lll>{\hsize=2\hsize}X}
\toprule
\textbf{Option} & \textbf{Défaut} & \textbf{Effet} & \textbf{Notes} \\
\midrule
\texttt{CMAKE\_BUILD\_TYPE} & Release & \texttt{-O3 -march=native} & Debug : \texttt{-O0 -g3} \\
\texttt{USE\_OMP} & ON & wfc\_omp\_lib + tests OMP & \\
\texttt{USE\_KOKKOS} & OFF & wfc\_kokkos\_lib + tests Kokkos & \texttt{Kokkos\_ROOT} requis \\
\texttt{USE\_LTO} & OFF & Link-Time Optimization & +3.4\% mesuré, A/B sur 11 runs \\
\texttt{BUILD\_DUNGEON} & OFF & wfc\_dungeon CLI + UE5 plugin & \\
\texttt{USE\_TSAN} & OFF & ThreadSanitizer (race detection) & Incompatible ASAN, LTO \\
\texttt{USE\_ASAN} & OFF & AddressSanitizer + UBSan & Incompatible LTO \\
\texttt{BUILD\_TESTING} & ON & Suite de tests ctest & \\
\bottomrule
\end{tabularx}
\caption{Options CMake exposées par le projet.}
\label{tab:cmake_options}
\end{table}

% =============================================================================
\section{Conventions de code}
\label{sec:impl:conventions}
% =============================================================================

\begin{itemize}
    \item C++17, sans dépendance externe runtime à part \texttt{stb\_image\_write.h} (third\_party, public domain).
    \item \texttt{Value = std::uint8\_t} partout (alphabet jusqu'à 256 valeurs).
    \item Headers publics dans \texttt{include/wfc/}, sources dans \texttt{src/}. Headers internes dans \texttt{include/wfc/internal/}.
    \item Aucun warning sous \texttt{-Wall -Wextra -Wpedantic -Werror} (g++ 14+, clang).
    \item Pas de \texttt{std::cout}/\texttt{printf} dans \texttt{src/} : I/O réservée aux apps.
    \item Pas de framework de test externe : \texttt{tests/test\_helpers.hpp} (40 lignes) suffit pour \texttt{WFC\_CHECK} / \texttt{WFC\_CHECK\_EQ} avec compteur global.
\end{itemize}

% =============================================================================
\section{Plates-formes vérifiées}
\label{sec:impl:plateformes}
% =============================================================================

\begin{table}[H]
\centering
\begin{tabularx}{\textwidth}{lXX}
\toprule
\textbf{Plateforme} & \textbf{Toolchain} & \textbf{Notes} \\
\midrule
Ubuntu 24.04 / WSL2 & g++ 13.3, OpenMP 4.5 & Machine de développement principale \\
Windows 11 (MinGW UCRT) & MSYS2 + g++ 16.1, OpenMP 5.2, Ninja & 13/13 tests passent ; perf parallèle limitée par heap lock MinGW \\
Romeo (RHEL 9, EPYC 9654) & spack + gcc 14.2, OpenMP 4.5 & Plate-forme cible HPC, 192 cœurs, 8 NUMA \\
Romeo armgpu (Grace+H100) & spack + gcc 11.4 + nvcc 12.9 + Kokkos CUDA & Build et tests Kokkos GPU validés \\
\bottomrule
\end{tabularx}
\caption{Plates-formes utilisées pour le développement et les benchmarks.}
\label{tab:plateformes}
\end{table}
